% ==================================================
% Discrete Exponential
% Discrete Calculus --- Special Functions
% ==================================================

\subsection{Special Functions in Discrete Calculus}
\label{sec:special-functions}

\begin{exposition}
Just as continuous calculus has its distinguished special functions ---
the exponential \(e^{ax}\), the trigonometric functions, and the logarithm
--- discrete calculus has analogous functions that behave simply under the
difference operator.
\end{exposition}

\subsection*{The Discrete Exponential}

\begin{exposition}
In continuous calculus, the exponential function \(e^{ax}\) is its own
derivative (up to a constant):
\[
\frac{d}{dx} e^{ax} = a\, e^{ax}.
\]

The analogous function in discrete calculus is the \emph{discrete
exponential} \((1+a)^n\).
\end{exposition}

\begin{definition}[Discrete Exponential]
\label{def:discrete-exponential}
For a constant \(a \ne -1\) and \(n \in \mathbb{Z}\), the
\emph{discrete exponential} with base \((1+a)\) is the function
\[
E_a(n) = (1+a)^n.
\]
\end{definition}

\DefinitionalRoot

\begin{remark*}[Standard quantified statement]
\(\forall a \in \mathbb{R}\setminus\{-1\},\; \forall n \in \mathbb{Z}:\;
E_a(n) := (1+a)^n\).
\end{remark*}

\begin{remark*}[Predicate reading]
The displayed quantified statement is the predicate-level form of $label: its hypotheses name the input conditions, and its conclusion names the asserted structure or relation.
\end{remark*}

\begin{remark*}[Interpretation]
The discrete exponential is the geometric function whose ratio between
successive inputs is \(1+a\).
\end{remark*}

\begin{proposition}[Discrete exponential rule]
\label{prop:discrete-exponential-rule}
\[
\Delta\,(1+a)^n = a\,(1+a)^n.
\]
\hyperref[prf:discrete-exponential-rule]{\textit{Go to proof.}}
\end{proposition}

\begin{remark*}[Standard quantified statement]
\(\forall a \in \mathbb{R}\setminus\{-1\},\; \forall n \in \mathbb{Z}:\;
(1+a)^{n+1} - (1+a)^n = a(1+a)^n\).
\end{remark*}

\begin{remark*}[Predicate reading]
The displayed quantified statement is the predicate-level form of $label: its hypotheses name the input conditions, and its conclusion names the asserted structure or relation.
\end{remark*}

\begin{remark*}[Interpretation]
The forward difference of a discrete exponential is a scalar multiple of the
same discrete exponential.


The continuous exponential satisfies \(\tfrac{d}{dx}e^{ax} = ae^{ax}\).
The discrete exponential satisfies \(\Delta(1+a)^n = a(1+a)^n\). The
base \(e\) in continuous calculus is replaced by \((1+a)\) in discrete
calculus. When \(a\) is small, \((1+a) \approx e^a\), recovering the
continuous limit.


From the proposition, a discrete antiderivative of \(a(1+a)^n\) is
\((1+a)^n\) itself. Equivalently, a discrete antiderivative of \((1+a)^n\)
is \(\tfrac{(1+a)^n}{a}\). Combined with the Fundamental Theorem of Finite
Calculus, this gives the geometric series formula:
\[
\sum_{k=0}^{n-1}(1+a)^k
= \frac{(1+a)^n - 1}{a}
\qquad (a \ne 0).
\]
\end{remark*}


\begin{dependencies}
\begin{itemize}
  \item \hyperref[def:discrete-exponential]{Discrete Exponential}.
  \item \hyperref[def:forward-difference]{Forward Difference Operator}.
\end{itemize}
\end{dependencies}
